\begin{theorem}[Flat torus connections and commuting seams]%
    \label{thm:peps_torus_flat_connection_gauge}
    \lean{TNLean.PEPS.zmodTransport_step,
        TNLean.PEPS.IsTorusFlat.zmodTransport_rectangle,
        TNLean.PEPS.IsTorusFlat.commute_holonomies,
        TNLean.PEPS.torusTreeGauge,
        TNLean.PEPS.IsTorusFlat.torusTreeGauge_right,
        TNLean.PEPS.IsTorusFlat.torusTreeGauge_up,
        TNLean.PEPS.torusNativeRightTransport,
        TNLean.PEPS.torusNativeUpTransport,
        TNLean.PEPS.torusNativeRightTransport_closure,
        TNLean.PEPS.torusNativeUpTransport_closure,
        TNLean.PEPS.torusEdgeAssignment_ext,
        TNLean.PEPS.exists_regularVertexGauge_eq_torusClosure}
    \leanok
    \uses{lem:peps_torus_parallel_section,
        thm:peps_regular_torus_closure_cut_identification}
    Let $G$ be a group and let $a_{x,y},b_{x,y}\in G$ be rightward
    and upward transports on $\mathbb Z_w\times\mathbb Z_h$, where
    $w,h>0$. Suppose the connection is flat:
    $b_{x+1,y}a_{x,y}=a_{x,y+1}b_{x,y}$.
    Write $A_m=a_{m-1,0}\cdots a_{0,0}$ and
    $B_{x,n}=b_{x,n-1}\cdots b_{x,0}$, with empty products equal to
    the identity. The holonomies $A_w$ and $B_{0,h}$ commute.
    For the least nonnegative representatives of $x,y$, set
    $k_{x,y}=B_{x,y}A_x$. Then
    \begin{align}
        k_{x+1,y}^{-1}a_{x,y}k_{x,y}
         & =\begin{cases}A_w&x+1=0,\\1&x+1\ne0,\end{cases}
        \label{eq:peps_flat_horizontal_seam}                   \\
        k_{x,y+1}^{-1}b_{x,y}k_{x,y}
         & =\begin{cases}B_{0,h}&y+1=0,\\1&y+1\ne0.\end{cases}
        \label{eq:peps_flat_vertical_seam}
    \end{align}
    If $w,h\ge3$ and the transports come from operators $u_e$ on
    the ordered edges of the torus graph, the same gauge transforms
    $u_e$ into the closure assignment with $g=B_{0,h}^{-1}$ and
    horizontal closure element $A_w$. No trivial-holonomy
    assumption is made. This is the group-valued seam deformation
    of \cite[Section~5]{Schuch2010PEPS}.
\end{theorem}

\begin{proof}\leanok
    Induction across a rectangle gives
    \begin{align}
        B_{m,n}A_m
         & =(a_{m-1,n}\cdots a_{0,n})B_{0,n}.
        \label{eq:peps_flat_rectangle_transport}
    \end{align}
    Taking $m=w,n=h$ in
    (\ref{eq:peps_flat_rectangle_transport}) gives
    $B_{0,h}A_w=A_wB_{0,h}$. Before a coordinate wraps, the
    transport to its successor adds one factor on the left;
    at the wrap it leaves the full cycle holonomy.
    Applying these two cases to $k$ proves
    (\ref{eq:peps_flat_horizontal_seam}) and
    (\ref{eq:peps_flat_vertical_seam}).
    For the left regular matrices $L_q$, the change on an ordered
    edge $e:s\to t$ is the following contraction, with matrix products
    read from right to left:
    % Formula: L_{k_t^{-1}} L_{u_e} L_{k_s}=L_{k_t^{-1}u_e k_s}.
    % Source: SCP10 paper_v3.tex:1623--1655, oriented string deformation;
    % TorusFlatConnectionGauge.lean, exists_regularVertexGauge_eq_torusClosure.
    % Ink-to-index: L_q(beta,alpha)=1_{beta=q alpha}; the two internal
    % group indices are summed, beta is the output, alpha the input.
    % Boundary signature: (V,P)=(2,0) on both sides. No physical site is drawn.

    \begin{tenkzequation}
        \tenkzkernel
        \begin{tenkz}[rows={wire}, cols=3, bonds=none]
            \tn[at=(1,1), name=flatGauge0,
                ports={180:virtual:$\beta$,0:virtual}]{$L_{k_t^{-1}}$}
            \tn[at=(1,2), name=flatGauge1,
                ports={180:virtual,0:virtual}]{$L_{u_e}$}
            \tn[at=(1,3), name=flatGauge2,
                ports={180:virtual,0:virtual:$\alpha$}]{$L_{k_s}$}
            \tnwire{flatGauge0.0}{flatGauge1.180}
            \tnwire{flatGauge1.0}{flatGauge2.180}
        \end{tenkz}
        $=$
        \begin{tenkz}[rows={wire}, cols=1, bonds=none]
            \tn[at=(1,1), name=flatGauged0,
                ports={180:virtual:$\beta$,0:virtual:$\alpha$}]{$L_{k_t^{-1}u_e k_s}$}
        \end{tenkz}
    \end{tenkzequation}
    On an ordered seam edge the native orientation is reversed.
    Thus its ordered horizontal operator is $A_w^{-1}$ and its
    ordered vertical operator is $B_{0,h}^{-1}$, exactly the
    claimed closure assignment.
\end{proof}

\begin{theorem}[Flat insertions give commuting closure states]%
    \label{thm:peps_flat_insertion_closure_state}
    \lean{TNLean.PEPS.exists_sameState_torusClosure_of_isTorusFlat,
        TNLean.PEPS.exists_stateCoeff_eq_torusGClosure_of_isTorusFlat}
    \leanok
    \uses{thm:peps_torus_flat_connection_gauge,
        thm:peps_regular_closed_gauge,
        thm:peps_regular_torus_closure_cut_identification}
    Let $G$ be finite and let the torus periods be at least three.
    Suppose the site coefficients are invariant under simultaneous
    left translation of their incident group labels. Every flat
    inserted-bond network has the same physical state as a network
    with two commuting closure seams. For a single four-leg tensor,
    this is the actual closure vector $\psi_{g,h}$ of
    \cite[Section~5]{Schuch2010PEPS}. Neither injectivity nor
    isometry is needed for this equality.
\end{theorem}

\begin{proof}\leanok
    Apply Theorem~\ref{thm:peps_torus_flat_connection_gauge}.
    Vertex gauges leave the closed contraction unchanged by
    Theorem~\ref{thm:peps_regular_closed_gauge}.
    Theorem~\ref{thm:peps_regular_torus_closure_cut_identification}
    identifies the resulting ordered-edge contraction with the
    four-leg closure vector.
\end{proof}
